% Auto-generated by src/make_tables.py -- do not edit by hand.
\begin{table}[t]
\centering
\caption{The grid timings, separated into the scheme and the implementation. The same solve --- same grid, same stencil, same step count, checked to agree on the solved field to $1.3\times10^{-4}$ relative (the edge closure; the interior is identical) --- is run by the NumPy solver used everywhere else in this paper and by a fused parallel kernel (numba; the environment has no JAX). The implementation factor grows with the problem, from $2.5\times$ at $p=1$ to $11.9\times$ at $p=4$, and the achieved rate rises --- $p=1$: $0.01\to0.01$, $p=2$: $0.14\to0.53$, $p=3$: $0.58\to4.41$, $p=4$: $0.65\to7.65$, $p=5$: $0.67\to6.43$ GFLOP/s --- which is what the FLOP model in the text was there to correct for: the NumPy figure was measuring per-step array overhead as much as arithmetic. The consequence for the comparison is at $p=5$, where the NumPy grid was slower than the neural solver ($865$\,s against $226$\,s) and the compiled one is faster ($59$\,s). Both columns are the time-stepping sweeps only, at the same forced step count.}
\label{tab:compiled}
\scriptsize
\setlength{\tabcolsep}{4pt}
\begin{tabular}{rrrrrr}
\toprule
$p$ & $N$ & steps & NumPy (s) & kernel (s) & factor \\
\midrule
1 & 21 & 420 & 0.02 & 0.009 & $2.5\times$ \\
2 & 21 & 420 & 0.05 & 0.015 & $3.7\times$ \\
3 & 21 & 420 & 0.58 & 0.076 & $7.7\times$ \\
4 & 21 & 420 & 18.48 & 1.558 & $11.9\times$ \\
5 & 21 & 420 & 567.10 & 59.252 & $9.6\times$ \\
\bottomrule
\end{tabular}
\end{table}
